Let $(R, \mathfrak m)$ be a Noetherian local ring.
From the above it is clear that $\mathfrak m$ cannot be
generated by fewer than $\dim(R)$ variables.
By Nakayama's Lemma \ref{lemma-NAK} the minimal number
of generators of $\mathfrak m$ equals $\dim_{\kappa(\mathfrak m)}
\mathfrak m/\mathfrak m^2$. Hence we have the following
fundamental inequality
$$
\dim(R) \leq \dim_{\kappa(\mathfrak m)} \mathfrak m/\mathfrak m^2.
$$
It turns out that the rings where equality holds
have a lot of good properties. They are called
regular local rings.

